
\section{Photoionization Viability Assessment}


\hypertarget{species-concentration-and-conductivity-calculations}{%
\subsection{Species concentration and Conductivity calculations}\label{species-concentration-and-conductivity-calculations}}

We begin by modeling the combustion of seeded oxy-kerosene mixture with the Cantera python package. The input mass flow rates of $K_2CO_3$, $C_{12}H_{26}$, and $O_2$ are determined to achieve a specified equivalence ratio and potassium mass fraction. This mixture is initialized at a inlet temperature (300 K) then allowed to equilibrate at a fixed enthalpy and combustion chamber pressure (1e6 Pa).

We then model different potential gas areas of the MHD channel by equilibrating this gas mixture at various fixed temperature and pressures. Our Cantera model does not consider seed particle dissociation or formation. Figure 1 shows the calculated concentrations of the 10 highest-concentration species for our modeled oxy-fuel combustion gas mixture as function of temperature and atmospheric pressure.

We calculate the conductivity of our working fluids following the method outlined previously by Bedick et al.{[}14{]} Briefly, the collisional cross sections for alkali vapors in flames. Then the electron mobility, \(\mu_{e}\), is calculated using the species-dependent collisional cross sections of constituent gasses. Then the electrical conductivity is calculated with the free electron density, \(n_{e}\):

\begin{equation}
\sigma = e\mu_{e}n_{e}
\end{equation}

\includegraphics[width=4.09929in,height=3.95663in]{\repodir/final/figures/SI/output/Cantera_T.png}

Figure 1: a) Mole fractions of species with mole fraction of greater than 1e-3. The calculations are performed for oxy-kerosene combustion with \(\phi = 0.8\), and \(K_{wt} = 1\%\). b) Ionization fraction = {[}K+{]}/({[}K{]} + {[}K+{]}) and c) electrical conductivity.

\hypertarget{nonequilibrium-conductivity}{%
\subsection{Nonequilibrium conductivity}\label{nonequilibrium-conductivity}}

To study the fundamental potential of NEI, we develop a 0D model of an infinitesimal control volume of gas described by the thermal equilibrium properties computed above. NEI is modeled as a perturbation on the seeded gas mixture in chemical equilibrium. We consider an arbitrary ionization mechanism that requires a power input per unit volume, \(P_{in}\), to generate a nonequilibrium generation rate \(G_{NE}\) . We can then define an ionization efficiency, \(\eta\) in the following way:

\begin{equation}
{G}_{NE} = \frac{\eta P_{In}}{eE_{IP}}
\end{equation}


This equation is agnostic to the specific ionization mechanism, which would be encapsulated in \(\eta\). In this interpretation, \(G_{NE}\) represents the rate that electrons are lifted from the potential well of depth \(E_{IP}\). When \(\eta = 1\), the ionization process is perfectly efficient and for \(\eta < 1\) a fraction \(\eta - 1\) of the input energy is wasted. We note that \(\eta > 1\) is possible, particularly in the case of laser enhanced ionization where a laser excites an atom to a higher bound state, and a subsequent collisional excitation ionizes the atom. \(\eta\ \)on the order of 10 has been observed for metals in flames.{[}15{]} In this work, we consider $\eta =1$ or lower. 

We consider the control volume of this gas subjected to a steady $G_{NE}$ and calculate the power dependent electron concentration, \(n_{e}(P_{in}),\), assuming the system is in quasi-equilibrium. The total generation rate of electrons (G) equals the recombination rate (R).
\begin{equation}
G = R
\end{equation}

G will be given both by $G_{NE}$ as well as a thermal generation rate from potassium ionization, \(G_{th} = k_{th}n_{K}\), where $k_{th}$ is the thermal generation rate constant.

\begin{equation}
G = G_{th} + G_{NE} = k_{th}n_{K} + \ \frac{\eta P_{In}}{eE_{IP}}
\end{equation}

$k_{th}$ is determined from the Saha equation and the bimolecular recombination coefficient for potassium, $k_{r,b,K+}$. 

\begin{equation}
    k_{th} = 2.42e^{15}*T^{3/2}*e^{-E_{ion}/k_bT}*k_{r,b,K+}
\end{equation}


The recombination rate is given by

\begin{equation}
R = k_{r, b, K+}*n_{e}^2  + k_{r, b, O2}*n_e*n_{O2}
\end{equation}

Where we have assumed \(n_{K^{+}}\)= \(n_{e}.\ \)



Combining these expressions

\begin{equation}
0 = k_{r, b, K+}*n_{e}^2  + k_{r, b, O2}*n_e*n_{O2} - G
\end{equation}

solving for $n_{e}$

\begin{equation}
n_e = \frac{-k_{r, b, O2}*n_{O2} + \sqrt{k_{r, b, O2}^2*n_{O2}^2 + 4*k_{r, b, K+}*G}}{2*k_{r, b, K+}}
\end{equation}


\begin{equation}
n_e = \frac{-k_{r, b, O2}*n_{O2}}{2*k_{r,b,K+}} + \sqrt{\frac{k_{r, b, O2}^2*n_{O2}^2 + 4*k_{r, b, K+}*G}{4*k_{r,b,K+}^2}}
\end{equation}


\begin{equation}
n_e = \frac{-k_{r, b, O2}*n_{O2}}{2*k_{r,b,K+}} + 
\sqrt{(\frac{k_{r, b, O2}*n_{O2}}{4*k_{r,b,K+}})^2 + \frac{G}{k_{r,b,K+}}}
\end{equation}


In the regime that K+ recombination dominates

\begin{equation}
n_{e}(P_{in}) = \sqrt{\frac{{G_{th} + \ G}_{NE}(P_{in})}{k_{r}} = \ }\sqrt{\frac{k_{th}n_{K} + \frac{\eta P_{In}}{eE_{IP}}}{k_{r}}\ }
\end{equation}

conversely when the O2 recombination dominates

\begin{equation}
  n_e = \frac{G}{k_{r,b,O2}*n_{O2}}
\end{equation}

\hypertarget{figure-of-merit}{%
\subsection{Figure of merit}\label{figure-of-merit}}

We are interested in finding the regimes in which a power input \(dP_{in}\) will generate a correspondingly larger amount of MHD power output \(dP_{MHD}\). We therefore define the figure of merit,

\begin{equation}
\alpha = \frac{{dP}_{MHD}}{dP_{in}} = \frac{{dP}_{MHD}}{dn_{e}}\frac{dn_{e}}{dP_{in}}\ 
\end{equation}


and want to find the regimes in which \(\alpha > 1,\ \)We assume that \(\mu_{e}\) independent of \(n_{e}\) to arrive at

\begin{equation}
\alpha = K(1 - K)u^{2}B^{2}e\mu_{e}\frac{dn_{e}}{dP_{in}}
\end{equation}


Calculating \(\frac{dn_{e}}{dP_{in}}\) with equation 7 gives

\begin{equation}
\alpha\ (P_{in}) = \eta\frac{K(1 - K)u^{2}B^{2}e\mu_{e}}{2E_{IP}k_{r}n_{e}(P_{in})}
\end{equation}

TODO: Extra e on top here? should cancel?

This expression can be interpreted as the ratio of the increases in MHD power to recombination losses (\(P_{R} = E_{IP}R = \ E_{IP}k_{r}n_{e}\)) with an induced increase in \(n_{e}\). Equation 10 and equation 7 show that the efficiency decreases as a function of \(P_{in}\) and recombination losses increase. We chose to examine the best-case scenario where \(P_{in}\)= 0, represents the absolute best infinitesimal energy return on energy investment possible for a given set of parameters. The regimes where \(\alpha\ \left( P_{in} = 0 \right)\ \)is high do not necessarily correspond to regimes in which NEI will result in a large \(P_{MHD}\) and can be large when the conductivity (and recombination) is low and correspondingly actual MHD power density (\(P_{MHD}\)) is near low.

\hypertarget{photoionization}{%
\subsection{Photoionization}\label{photoionization}}

We apply our model to the specific case of UV photoionization. Photoionization, where potassium atoms are ionized by UV photons with energy above the potassium ionization potential, is a potential method of increasing the electrical conductivity of seeded combustion products. Photoionization is particularly attractive for contact targeting, due to the ability to direct and spatially confine light. Photoionization was considered as a method of increasing \(\sigma\) in MHD power systems early in MHD power research.{[}4{]} In that work, the potential of photoionization as an effective method of achieving sufficiently high electrical conductivity for MHD power was estimated. Those estimates were lacking in multiple ways:

\begin{itemize}
\item
  They considered blackbody radiation as the source of UV photons as UV laser technology was not developed yet. This is a highly inefficient way to generate UV photons. In addition, the use of lasers facilitates directional control of light, and advances in UV resistant fiber optics could potentially allow for precise application to specific areas.
\item
  The ionization of only a potassium vapor was considered, neglecting the absorption of light by other constituent gasses.
\item
  The calculations employed estimates of the potassium electron recombination rate, and experimental measurements of these reaction rates have been performed since then.
\end{itemize}

We have previously investigated photoionization in a HVOF jet using a 3D CFD model.{[}16{]}

We model photoionization by calculating the ionization efficiency for photoionization, \(\eta_{PI}\) (See equation 3). Considering a fraction of the incident photons that are absorbed by potassium, \(FA_{K}\), and that any photon energy above the ionization potential is wasted, we conclude

\begin{equation}
\eta_{PI} = \frac{E_{IP}}{E_{ph}}FA_{K} = \frac{E_{IP}}{E_{ph}}\frac{\delta}{\delta_{K}\ }FA\ \ 
\end{equation}

Where \(FA_{K}\) has been expressed in terms of the optical attenuation length of each individual species \(\delta_{i}\) and the total optical attenuation length of the gas mixture, \(\delta\), and the fraction of incident light absorbed by the entire gas mixture, FA. To understand the fundamental limit of photoionization we assume FA = 1. The case of FA for a practical system of an optical cavity with finite reflectivity mirrors is explored in the appendix.

\(\delta\) is calculated with the following equation:

\begin{equation}
\delta = \sum_{i}^{}\delta_{i} = \sum_{i}^{}{n_{i}Q_{i}}
\end{equation}

Where \(Q_{i}\) is the optical absorption cross section of the \emph{i}\textsuperscript{th} species. In this way \(\delta_{K} = n_{K}Q_{K}\)where the \(K\) stands for potassium.

To calculate \(\frac{\delta}{\delta_{K}\ }\) we utilize the compiled UV absorption cross section data outlined in \ref{fig:SI_UV_abscs} on the UV optical absorption cross sections for the major species in the calculated gas mixtures.

Figure 6 shows\(\ \delta\), \(\delta_{K}\), and \(\frac{\delta}{\delta_{K}\ }\)for various UV wavelengths and temperatures.

\includegraphics[]{\repodir/modeling/viability/figure_panels/output/atten_length.png}

In the following work we analyze these gasses excited by 248 nm light. This is primarily as that is the longest wavelength of an excimer laser (KrF) that can still ionize potassium. However, the fraction of light absorbed by potassium is not particularly dependent on the wavelength, as can be seen in Figure 6. Although excimer lasers are pulsed we choose to analyze steady state illumination here. This is because for a given time-averaged laser intensity the time-averaged electron density will be higher with a steady state light source. The higher peak intensities of pulsed laser sources will create higher peak electron densities and a correspond higher rate of second order recombination. Therefore, considering steady state illumination corresponds to an upper bound on the electron density achievable with a given photoionization energy input. This assumption also allows for the calculation of equilibrium electron density by allowing for the equivalence of recombination rate and generation rate.

\clearpage

\subsection{Supplemental PI Viability Results}

\begin{figure}
\centering
\includegraphics[width=\textwidth]{\repodir/modeling/viability/figure_panels/output/beta_boundary_layer.png}
\caption{Supplemental results}
\end{figure}
